% © 2026 Ing. David Jaroš — CC BY-NC-ND 4.0
% UBT-AI-PROVENANCE-BEGIN
% schema: ubt-ai-provenance/v1
% tier: B_machine_verified
% ai_assistance: disclosed
% human_review: machine-verification
% editorial_responsibility: Ing. David Jaroš
% policy: ../../AI_PROVENANCE.md
% notice: Machine-verified against named sources or verifiers; individual attestation is not claimed.
% UBT-AI-PROVENANCE-END

\documentclass[11pt,a4paper]{article}
\usepackage[margin=2.5cm]{geometry}
\usepackage{amsmath,amssymb,amsthm,mathtools,booktabs}
\usepackage[most]{tcolorbox}
\newtheorem{theorem}{Theorem}
\newtheorem{corollary}[theorem]{Corollary}
\newtheorem{proposition}[theorem]{Proposition}
\newtheorem{remark}[theorem]{Remark}
\newcommand{\Tr}{\operatorname{Tr}}
\newcommand{\Ei}{\operatorname{E_1}}
\title{GAP-10D Endgame Audit:\\
Axiomatic Underdetermination and the Correct Induced-Einstein Coefficient}
\author{Ing. David Jaro\v{s}}
\date{26 July 2026}
\begin{document}\maketitle

\begin{abstract}
We determine exactly how far the present UBT axioms can derive Einstein
dynamics.  First, we prove an underdetermination theorem: the covariant-tetrad
map, central metric identity, split-jet right inverse, and all currently stated
local symmetries do not fix the coefficient of a curvature action.  Infinitely
many actions with different Newton constants share the same UBT kinematics.
Thus an unconditional derivation of $G$ from the existing kinematic axioms is
logically impossible.  Second, we compute the curvature term generated by a
proper-time regulated one-loop determinant of a specified Laplace-type
$\Theta$ Hessian on $M^4\times S^1_\psi$.  The result gives an exact conditional
formula for the induced Newton coefficient, including the full Kaluza--Klein
heat trace.  At the self-dual cutoff $\Lambda R_\psi=1$ the dimensionless KK
factor is $1.303410251859\ldots$.  This closes the form of the induced-Einstein
bridge, but the physical mode count, nonminimal curvature coupling, and cutoff
identification must be derived or explicitly adopted before the numerical
Newton constant is a UBT prediction.
\end{abstract}

\section{What the current axioms do not determine}
Let $g[\Theta]$ denote any metric admitted by the local split-jet construction.
The current geometric axioms specify the field space and the map to the metric,
but not a unique scalar functional of that metric.

\begin{theorem}[Curvature-coefficient underdetermination]\label{thm:under}
For every real constant $c$, the action family
\begin{equation}
 S_c[\Theta]=S_{\rm locked}[\Theta]
 +c\int d^4x\,\sqrt{-g[\Theta]}\,R[g[\Theta]]
 \label{eq:family}
\end{equation}
has the same covariant-tetrad definition, central metric identity, local
Lorentz symmetry, split-jet representability, and kinematic connection
reconstruction.  Distinct values of $c$ give distinct Newton constants.
Therefore those kinematic statements cannot determine $c$.
\end{theorem}

\begin{proof}
Every item listed in the theorem is a definition or an algebraic/differential
identity on the field space and is independent of the coefficient multiplying
a separately invariant scalar functional.  The second term in
Eq.~\eqref{eq:family} is diffeomorphism and local-Lorentz invariant for every
$c$.  Changing $c$ leaves all kinematic identities unchanged but changes the
metric Euler--Lagrange equation and the Newton coefficient.  Hence $c$ is not a
logical consequence of the current kinematic axioms.
\end{proof}

The same argument applies to $R^2$, $R_{\mu\nu}R^{\mu\nu}$, and other higher-
derivative invariants.  Lovelock uniqueness identifies the two-derivative
infrared form once its assumptions are adopted; it does not determine its
overall coefficient or derive the action from $\Theta$.

\begin{corollary}[Absolute-scale obstruction]\label{cor:scale}
The present normalization convention treats $\mathcal N_0$ as a global
unit-setting constant.  Under a change of mass unit by a positive factor $a$,
\begin{equation}
 \mathcal N_0\mapsto a^4\mathcal N_0,\qquad
 \Lambda\mapsto a\Lambda,\qquad
 G\mapsto a^{-2}G,
\end{equation}
while all dimensionless geometric identities are unchanged.  Consequently an
absolute numerical value of the dimensionful constant $G$ cannot follow from
the scale-free kinematic identities alone.  The meaningful candidate
prediction is a dimensionless ratio such as $G\sqrt{\mathcal N_0}$; even that
ratio requires the physical Hessian, measure and matching prescription.
\end{corollary}

\section{Specified one-loop branch}
A genuine induced-gravity calculation must begin with a specified quadratic
fluctuation operator, not merely an abstract Dirac operator; see the standard
heat-kernel review~\cite{vassilevich2003} and the induced-gravity proposal of
Sakharov~\cite{sakharov1967}.  Assume that after
background expansion, gauge fixing, and removal of nonphysical directions the
Euclidean Hessian contains $N_B$ real bosonic modes with Kaluza--Klein label
$n\in\mathbb Z$ and Laplace-type operators
\begin{equation}
 P_n=-\nabla^2+m^2+\frac{n^2}{R_\psi^2}+\xi R+\mathcal E_0,
 \label{eq:Pn}
\end{equation}
where $\mathcal E_0$ has no term proportional to the scalar curvature.  The
one-loop contribution is
\begin{equation}
 \Gamma_1=\frac12\sum_n\Tr\log P_n
 =-\frac12\int_{\Lambda^{-2}}^\infty\frac{ds}{s}
 \sum_n\Tr e^{-sP_n},
 \label{eq:proper}
\end{equation}
using a proper-time cutoff $\Lambda$.

For the four-dimensional heat kernel,
\begin{equation}
 \Tr_{M^4}e^{-sP_n}
 =\frac{N_B}{(4\pi s)^2}\int d^4x\sqrt g\,
 e^{-s(m^2+n^2/R_\psi^2)}
 \left[1+s\left(\frac16-\xi\right)R+O(s^2)\right].
 \label{eq:hk}
\end{equation}
Define the exact regulated KK integral
\begin{equation}
 \mathcal I_1(\Lambda,R_\psi,m)
 :=\int_{\Lambda^{-2}}^\infty ds\,s^{-2}e^{-m^2s}
 \vartheta_3\!\left(0,e^{-s/R_\psi^2}\right).
 \label{eq:I1}
\end{equation}

\begin{theorem}[Induced Einstein coefficient]\label{thm:induced}
Under the stated Hessian and regulator assumptions, the curvature part of the
one-loop Euclidean action is
\begin{equation}
 \boxed{
 \Gamma_1\big|_R
 =-\frac{N_B(1-6\xi)}{192\pi^2}\,
 \mathcal I_1(\Lambda,R_\psi,m)
 \int d^4x\sqrt g\,R.}
 \label{eq:GammaR}
\end{equation}
Matching the Euclidean convention
$S_E=-\frac1{16\pi G}\int\sqrt g\,R+\cdots$ gives
\begin{equation}
 \boxed{
 \frac1{G_{\rm ind}}
 =\frac{N_B(1-6\xi)}{12\pi}\,
 \mathcal I_1(\Lambda,R_\psi,m).}
 \label{eq:Gind}
\end{equation}
For minimally coupled bosonic modes, $\xi=0$, the induced Newton constant has
the physical sign.
\end{theorem}

\begin{proof}
Insert Eq.~\eqref{eq:hk} into Eq.~\eqref{eq:proper}.  The coefficient of $R$
is
\[-\tfrac12(4\pi)^{-2}N_B(\tfrac16-\xi)\mathcal I_1,\]
which is Eq.~\eqref{eq:GammaR}.  Matching coefficients yields
Eq.~\eqref{eq:Gind}.
\end{proof}

\section{Low-energy and self-dual limits}
If $\Lambda R_\psi\ll1$ and $m\ll\Lambda$, nonzero KK modes are exponentially
suppressed and
\begin{equation}
 \mathcal I_1=\Lambda^2+O(e^{-1/(\Lambda R_\psi)^2}).
\end{equation}
At the self-dual identification $\Lambda R_\psi=1$ and $m=0$, put
$s=u/\Lambda^2$.  Then
\begin{equation}
 \mathcal I_1=C_\psi\Lambda^2,
 \qquad
 C_\psi:=\int_1^\infty\frac{du}{u^2}
 \vartheta_3(0,e^{-u}).
\end{equation}
Termwise integration gives the rapidly convergent exact representation
\begin{equation}
 C_\psi=1+2\sum_{n=1}^\infty
 \left[e^{-n^2}-n^2E_1(n^2)\right]
 =1.303410251859279308\ldots .
 \label{eq:Cpsi}
\end{equation}
Therefore
\begin{equation}
 \frac1{G_{\rm ind}}
 =\frac{N_B(1-6\xi)C_\psi}{12\pi}\Lambda^2.
 \label{eq:Gself}
\end{equation}

Dimensional analysis of $E_\mu=\mathcal N_0^{-1/2}D_\mu\Theta$ with
$[\Theta]=1$ gives $[\mathcal N_0]=4$.  If one explicitly adopts the
scale-identification axiom
\begin{equation}
 \Lambda=\mathcal N_0^{1/4}=R_\psi^{-1},
 \label{eq:scaleaxiom}
\end{equation}
then the reduced Planck mass obeys
\begin{equation}
 \boxed{
 \bar M_{\rm Pl}^2:=\frac1{8\pi G_{\rm ind}}
 =\frac{N_B(1-6\xi)C_\psi}{96\pi^2}\sqrt{\mathcal N_0}.}
 \label{eq:Mpl}
\end{equation}
For the unconstrained eight-real-component biquaternion count and $\xi=0$,
$\bar M_{\rm Pl}=0.10490594\,\mathcal N_0^{1/4}$.  This numerical relation is
conditional because the physical gauge-fixed mode count need not equal eight.

\section{Why the earlier five-dimensional note is not a closure proof}
The exploratory May 2026 file
\texttt{research\_tracks/T3\_ALPHA/seeley\_dewitt\_coefficients.tex} cannot be
used as a GR coefficient derivation:
\begin{enumerate}
 \item it introduces a spectral Dirac operator that is not shown to be the
 quadratic Hessian of the canonical $\Theta$ action;
 \item it double-counts the Lichnerowicz $R/4$ term by writing it once inside
 $D_{M^4}^2$ and once again in the product operator;
 \item it mixes real and complex fibre dimensions in the trace count;
 \item a five-dimensional local $a_2$ coefficient and a four-dimensional
 low-energy Newton coefficient have different cutoff scaling;
 \item the finite KK theta factor belongs to the regulated full heat trace
 Eq.~\eqref{eq:I1}, not to a local Seeley--DeWitt coefficient by itself.
\end{enumerate}
The corrected computation is Eq.~\eqref{eq:Gind}.

\section{Closure boundary}
The induced formula is exact once $(P_n,N_B,\xi,\Lambda)$ are specified.  The
current repository does not yet derive all four from the finalized constrained
$\Theta$ measure.  In particular, GAP-Q still lists the functional measure,
gauge/ghost quotient, and UV scale as open.  Therefore one may adopt a
well-defined \emph{spectral-induced effective branch}, but one may not relabel
it as an unconditional consequence of the earlier kinematic axioms.

\begin{tcolorbox}[colback=green!6!white,colframe=green!55!black,title=Exact status]
\textbf{GAP-10D-UNDERDETERMINATION: CLOSED AS NO-GO [L1].}
The existing geometric axioms do not determine a curvature coefficient or
Newton constant; an additional dynamical/quantum principle is necessary.
Corollary~\ref{cor:scale} separately shows that the absolute dimensionful
scale cannot be predicted before the physical scale represented by
$\mathcal N_0$ is fixed.

\medskip
\textbf{GAP-10D-A2-FORM: CLOSED CONDITIONALLY [L1].}
For a specified gauge-fixed Laplace-type $\Theta$ Hessian, proper-time one-loop
integration generates the Einstein--Hilbert term with the exact coefficient
Eq.~\eqref{eq:Gind}, including the full compact-$\psi$ KK trace.

\medskip
\textbf{GAP-10D-SPECTRAL-IR: CLOSED CONDITIONALLY [L1].}
If the Hessian, physical mode count, curvature coupling, and cutoff/scale
identification are adopted or independently derived, the two-derivative
infrared metric action is Einstein--Lambda, with higher-curvature terms
starting at the next heat-kernel order.

\medskip
\textbf{GAP-10D: NARROWED, not unconditionally closed.}
The remaining first-principles lemma is no longer the heat-kernel algebra.  It
is derivation of $(N_B,\xi,\Lambda)$ and the constrained fluctuation measure
from the finalized UBT action, without fitting $G$.
\end{tcolorbox}

\begin{thebibliography}{9}
\bibitem{sakharov1967}
A.~D.~Sakharov, ``Vacuum quantum fluctuations in curved space and the theory
of gravitation,'' \emph{Dokl. Akad. Nauk SSSR} \textbf{177}, 70--71 (1967)
[\emph{Sov. Phys. Dokl.} \textbf{12}, 1040--1041 (1968)].
\bibitem{vassilevich2003}
D.~V.~Vassilevich, ``Heat kernel expansion: user's manual,''
\emph{Phys. Rept.} \textbf{388}, 279--360 (2003).
\end{thebibliography}
\end{document}
